\par\bigskip\Needspace{9\baselineskip}
\pdfbookmark[2]{Second Method, [2]}{ceres-second-2}
\begin{center}[2.]\end{center}
\begingroup\small
\noindent\hspace*{\parindent}[For calculating the perturbations in latitude, the unperturbed orbit of Ceres is to be taken as the fundamental plane. We must therefore first calculate the relative orientation of the orbits of {\AstroSymbols ⚳} and {\AstroSymbols ♃}\refstepcounter{footnote}\hyperlink{ceresIII2-note}{\textsuperscript{*)}}\begingroup
\makeatletter\long\def\@makefntext#1{\noindent[*) #1}\makeatother
\footnotetext{\hypertarget{ceresIII2-note}{}See Article 15 of Exposition d’une nouvelle méthode de calculer les perturbations planétaires, printed below, or the somewhat fuller account in the calculation of the perturbations of Pallas.]}
\endgroup{}.

If we take \(N\) and \(N'\) to be the longitudes of the ascending nodes of the two orbits and assume
\[
\begin{array}{@{}r@{\;}l@{\qquad}r@{\;}l@{}}
i=&10^{0}37'54''&i'=&1^{0}18'51''\\
N=&80\ 57\ 43&N'=&98\ 26\ 27,
\end{array}
\]
we obtain
\[
\text{Mutual inclination of the two orbits}=\mathrm{J}=9^{0}23'11'',\qquad \log\sin\mathrm{J}=9{,}2124316,
\]
and for the longitude of the ascending node of the {\AstroSymbols ⚳} orbit on the {\AstroSymbols ♃} orbit we have
\[
\begin{array}{@{}c@{\;}c@{\;}c@{\;}l@{\quad}l@{}}
\text{measured}&\text{in}&\text{the}&\text{{\AstroSymbols ⚳} orbit:}&78^{0}32'26''\\
\text{»}&\text{»}&\text{»}&\text{{\AstroSymbols ♃} orbit:}&\AscendingNode=78^{0}34'38''.
\end{array}
\]

If in equation \(*\) for \(\theta\) (p. 381) we replace the longitude \(v\) with the true anomaly \(V\), we obtain
\[
\theta=\frac{\mu a}{\BareRoot{(1-ee)}}
\left\{\begin{aligned}
&\cos V\!\int r\left(\frac{\mathrm{d}R}{\mathrm{d}z}\right)\sin V.n\,\mathrm{d}t\\
&-\sin V\!\int r\left(\frac{\mathrm{d}R}{\mathrm{d}z}\right)\cos V.n\,\mathrm{d}t
\end{aligned}\right\}.
\]
Now (cf. p. 378) we have
\[
\begin{gathered}
\left(\frac{\mathrm{d}R}{\mathrm{d}z}\right)=z'\left(\frac{1}{r'^{3}}-\frac{1}{\mathfrak{R}^{3}}\right)\mathpunct{\raisebox{.45ex}{,}}\qquad
\mathfrak{R}=\BareRoot{(rr+r'r'-2rr'\cos w)}.
\end{gathered}
\]

Let \(l\) and \(b\) denote the heliocentric longitude and latitude of Ceres, referred to the {\AstroSymbols ♃} orbit as the fundamental plane, and let \(v'\) denote the true longitude of {\AstroSymbols ♃} in its orbit. Then
\[
\cos w=\cos b\cos(v'-l).
\]

Set
\[
\begin{array}{@{}l@{\qquad}l@{}}
l=v+\delta&D=\text{\AstroSymbols ♃}-\text{\AstroSymbols ⚳}\\
v=\text{\AstroSymbols ⚳}+\epsilon&\alpha=\epsilon+\delta,\\
v'=\text{\AstroSymbols ♃}+\epsilon'&
\end{array}
\]
where \(v\) is the true longitude of {\AstroSymbols ⚳} in its orbit, \(\delta\) is the reduction of this longitude to the {\AstroSymbols ♃} orbit up to sign, {\AstroSymbols ⚳} and {\AstroSymbols ♃} are the mean longitudes, and \(\epsilon=\text{Eq. of centre {\AstroSymbols ⚳}}\) and \(\epsilon'=\text{Eq. of centre {\AstroSymbols ♃}}\). We then obtain
\[
\cos w=\cos b\cos(D+\epsilon'-\alpha).
\]
\SourceLeaf{409}Hence
\[
\frac{1}{\mathfrak{R}^{3}}=\frac{1}{(rr+r'r'-2rr'\cos b\cos(D+\epsilon'-\alpha))^{\frac{8}{2}}}\mathpunct{\cdot}
\]

We can expand this expression into the series
\[
\frac{2}{\mathfrak{R}^{3}}=\tfrac12 Q+Q'\cos(D+\epsilon'-\alpha)+Q''\cos2(D+\epsilon'-\alpha)+\ldots
\]
and write
\[
\begin{aligned}
1)\quad\frac{2}{\mathfrak{R}^{3}}-\frac{2}{r'^{3}}={}&\tfrac12 Q^{(0)}+Q'\cos(D+\epsilon'-\alpha)\\
&+Q''\cos2(D+\epsilon'-\alpha)+\ldots
\end{aligned}
\]
where the \(Q^{(i)}\) are functions of \(r\), \(r'\) and \(\cos b\), with \(Q^{(0)}=Q-\frac{4}{r'^{3}}\).

We also have
\[
z'=r'\sin\mathrm{J}\sin(v'-\AscendingNode)=r'\sin\mathrm{J}\sin(\text{\AstroSymbols ♃}+\epsilon'-\AscendingNode);
\]
so that the expression for \(\theta\) becomes
\[
\begin{aligned}
2)\enspace\theta={}&-\frac{\mu a\sin\mathrm{J}}{4\BareRoot{(1-ee)}}\Biggl\{\pm\genfrac{}{}{0pt}{}{\cos}{\sin}V\!\int rr'\genfrac{}{}{0pt}{}{\sin}{\cos}V.\\
&\left\{\begin{array}{@{}l@{\;}l@{}}
Q'\sin(\text{\AstroSymbols ⚳}+\alpha-\AscendingNode)&\\[3pt]
+Q^{(0)}\sin(D+\text{\AstroSymbols ⚳}+\epsilon'-\AscendingNode)&-Q''\sin(D-\text{\AstroSymbols ⚳}+\epsilon'-2\alpha+\AscendingNode)\\[3pt]
+Q'\sin(2D+\text{\AstroSymbols ⚳}+2\epsilon'-\alpha-\AscendingNode)&-Q'''\sin(2D-\text{\AstroSymbols ⚳}+2\epsilon'-3\alpha+\AscendingNode)\\[3pt]
+Q''\sin(3D+\text{\AstroSymbols ⚳}+3\epsilon'-2\alpha-\AscendingNode)&-Q^{\mathrm{V}}\sin(3D-\text{\AstroSymbols ⚳}+3\epsilon'-4\alpha+\AscendingNode)\\[3pt]
+\ldots&-\ldots
\end{array}\right\}n\,\mathrm{d}t
\end{aligned}
\]
where the abbreviations \(\displaystyle\pm\genfrac{}{}{0pt}{}{\cos}{\sin}\) and, correspondingly, \(\displaystyle\genfrac{}{}{0pt}{}{\sin}{\cos}\) are understood to stand for the sum of these \textls[180]{two} terms.]
\par\endgroup
\begin{center}\rule{.30\linewidth}{.4pt}\end{center}
